\section{Indigenous Education and Practices}\label{indigenous-education-and-practices}

\stanceinfo{\textbf{CELC 2024} [2024-01-05]}{2024-01-05}

\officialstance{The Canadian Federation of Engineering Students acknowledges the harm done to Indigenous Peoples and recognizes that true reconciliation is an ongoing process. The role the engineer holds within reconciliation includes the rebuilding of our relationship with Indigenous Peoples. To incorporate Indigenous perspective into engineering problem-solving, engineering curricula must include Indigenous knowledge, as well as institutions must provide additional efforts to increase Indigenous representation in engineering.}

\stanceheading{The Issue}

Indigenous issues, knowledge, and socially responsive reconciliation should be prioritized within engineering education. Indigenous Peoples have undergone hardships throughout Canadian history, including the dispossession of ancestral lands and being denied the opportunity to lead as a primary consultant for issues relevant to their heritage. An intersectional perspective is fundamental to the understanding and mutual respect that must be established in order to move towards positive change.

Decolonization strives to create a safe space for all Indigenous staff, faculty, and students at institutions; more importantly protects local Indigenous Peoples to instead promote collaboration and local teachings. Due to the lack of localized Indigenous integration, there exists a disparity between education on history and the traditional engineering curriculum. Localized Indigenous education teaches how to communicate effectively and respectfully with Indigenous stakeholders and includes when to engage and what issues need to be addressed specific to the geographic region.

Within the CFES, we believe Indigenous education and engineering should hold an emphasis on the interactions with Indigenous peoples. We believe that the Canadian Engineering Accreditation Board currently lacks localized Indigenous education as a Graduate Attribute, in order to observe an intersectional perspective and create well-rounded engineers that are able to translate Indigenous knowledge and communicate with Indigenous communities into their engineering profession. Changes must be made to educate individuals, including changing curriculum and having conversations over Indigenous practices and improvements within engineering education.

\stanceheading{The Students' Position}
\begin{itemize}
 \item There is a lack of Indigenous representation in engineering when we consider only engineers in their main working years (age 25-64) with an education at a bachelor level or above, only 0.93\% identify as Indigenous [10].
 \begin{itemize}
  \item In 2021, 49.2\% of Indigenous people in their working years had completed a postsecondary certificate, degree or diploma, a rate lower than that of non-Indigenous people at 68.0\% [11]
  \item In 2021, Indigenous Peoples accounted for 5.0\% of the total population in Canada, around 1.8 million Indigenous people [12]
 \end{itemize}
 \item Increasing representation of Indigenous people in post secondary institutions should be a top priority with the assistance of government support for post-secondary education.
 \item Providing increased representation in the engineering profession can respond to the Truth and Reconciliation Commission's calls to action.
 \item Indigenous knowledge and wisdom must be a fundamental part of engineering education, to contribute perspective, traditional knowledge, and sustainable practices.
 \item Engineers interact with and directly impact Indigenous communities through the development, maintenance, and modification of infrastructure and engineering projects, therefore an intercultural understanding based in mutual respect must be found to promote positive communication and contribute to Reconciliation.
\end{itemize}

\stanceheading{Recommendations to Partners, Stakeholders, and Other Entities}
\begin{itemize}
 \item The CFES calls on Engineers Canada to hear Indigenous engineering student voices for the improvement of representation of Indigenous Peoples.
 \item The CFES calls on Engineers Canada to involve the CFES in the DIEEN subnetwork to increase CFES knowledge for resources and surveys.
 \item The CFES calls the Canadian Engineering Accreditation Board (CEAB) to improve upon the current Graduate Attributes and expand to include local Indigenous Practices and Education to include communication with Indigenous communities and Indigenous knowledge and wisdom.
 \item The CFES calls on the Engineering Deans Canada to include Indigenous Education and Practices into strategic plans, to ensure institutions hold this practice in high regard when making decisions.
 \item The CFES calls on Engineering Deans Canada to integrate Indigenous Education and Practices into their institutions by way of activities, lectures and events following the example of other post-secondary institutions.
 \item The CFES calls on Engineering Deans Canada to promote the Indigenous recruitment within engineering, providing resources and making engineering accessible to Indigenous Peoples.
\end{itemize}

\stanceheading{What the CFES is doing}
\begin{itemize}
 \item CFES conferences currently feature Indigenous Education and Practice through speakers, workshops, and different activities.
 \item The CFES is involved in the iBET PhD Project that supports Canadian Indigenous and Black professors in engineering and computer science fields across Canada.
 \item The CFES uses the National Survey to create a demographic landscape, including the identification of Indigenous peoples and participation in engineering.
\end{itemize}

\stanceheading{What the CFES plans to do}
\begin{itemize}
 \item Connect with the Engineers Canada community of practice called Decolonizing and Indigenizing Engineering Education Network (DIEEN) to address Indigenous access to engineering education and truth and reconciliation within engineering education [13].
 \item Gather information from Engineering Deans Canada from institutions that have established post secondary engineering initiatives to promote Indigenous education to develop a framework for advocacy surrounding Indigenous Education (as per Engineers Canada documentation):
 \begin{itemize}
  \item University of Alberta Indigenous Peoples and Technoscience Online Course
  \item Indigenous Futures in Engineering at Queen's University
  \item University of Manitoba's Engineering Access Program
  \item University of Saskatchewan Indigenous Engineering Access Program
 \end{itemize}
 \item Work with other partner organizations such as the American Indian Science and Engineering Society (AISES) to gain further insight and partner for further advocacy efforts.
 \item Analyze current information and surveys that have been collected by DIEEN and Engineers Canada to create an initial parameter for data collection and qualitative/quantitative analysis to then integrate questions about Indigenous Education into the National Survey.
\end{itemize}

\newpage
